\documentclass{article}
\title{Practice Problem 2.33}
\author{CSAPP}
\usepackage{booktabs, parskip, amssymb, tabularx}

\begin{document}
    \maketitle
    \begin{flushleft}
        We can represent a bit pattern of length $w=4$ with a single $hex$ digit. For a 
        two's complement interpretation of these digits, fil in the following table two
        determine the additive inverses of the digits shown:
    \end{flushleft}
    \smallskip

    \begin{center}
        \begin{tabular}{cccc}
            \multicolumn{2}{c}{$x$} & \multicolumn{2}{c}{$-^t_4x$}\\
            \midrule
            \texttt{Hex} & \texttt{Decimal} & \texttt{Decimal} & \texttt{Hex}\\
            \midrule
            $2$ & $2$  & $-2$ &$E$\\
            $3$ & $3$  & $-3$ &$D$\\
            $9$ & $9$  & $-9$ &$7$\\
            $B$ & $11$ & $5$  &$5$\\
            $C$ & $12$ & $4$  &$4$\\
        \end{tabular}
        \\
        The bits pattern are the same as for unsigned negation
    \end{center}
\end{document}